% GENERATED FILE. Edit canonical lessons, then run npm run generate:books.
% canonical-source-hash: fnv1a64:7332bba17b3feef3
% canonical-lessons: FA-C488-bahush, FA-C488-shukh, FA-C488-khosh-akhlaq, FA-C488-sekhavatmand, FA-C488-delsuz

\chapter{Good Character}
\label{ch:fa-good-character}
\hlchaptermodality{488}

\begin{chapteropening}
\textbf{By the end of this chapter:} \emph{I can describe good qualities in people.}

Complete the last lesson of chapter 488: I can describe good qualities in people.
\end{chapteropening}

\section[\texorpdfstring{\fa{باهوش} (bâhush)}{باهوش (bâhush)}]{\fa{باهوش} (bâhush) — clever, intelligent}

\label{lesson:FA-C488-bahush}



\begin{quote}
Before the new one: say the Persian for self-confidence, then the Persian for courage.
\end{quote}

\subsection*{You'll want to know: \fa{باهوش}}
\textbf{\fa{باهوش}} (\emph{bâhush}) — \textquotedblleft{}clever, intelligent\textquotedblright{}.

\begin{quote}
\fa{دخترشان} \fa{خیلی} \fa{باهوش} \fa{است}.
\end{quote}

\begin{quote}
\emph{dokhtarashân kheyli bâhush ast.}
\end{quote}

\begin{quote}
Their daughter is very clever.
\end{quote}

\emph{bâ} 'with' + \emph{hush} 'intelligence, wits'. The opposite in speech is \emph{kamhush}.

\begin{cousinweb}
\emph{bâhush} is built on native Persian \emph{bâ}, 'with', and \emph{hush}, 'wits', which is usually traced to Middle Persian \emph{hōš}. EtymDB has no entry for it.
\end{cousinweb}

\begin{grammarlens}[title={What you've built}]
One more word for telling, explaining and describing at B1.
\end{grammarlens}

\subsection*{Your turn}
\begin{itemize}
\raggedright
  \item \emph{Say it:} \emph{bâhush}
  \item \emph{Say it:} \emph{bâhush}, once more
  \item \emph{Say it:} a sentence of your own with \emph{bâhush}
  \item \emph{Recall:} say the Persian for self-confidence, then the Persian for courage, then say \emph{bâhush} again
\end{itemize}

\begin{tcolorbox}[breakable,colback=teal!4,colframe=teal!35!black,title={Before you move on}]
What does \fa{باهوش} mean? (\textquotedblleft{}Clever, intelligent\textquotedblright{}.) Say it once more.
\end{tcolorbox}
\section[\texorpdfstring{\fa{شوخ} (shukh)}{شوخ (shukh)}]{\fa{شوخ} (shukh) — witty, playful}

\label{lesson:FA-C488-shukh}



\begin{quote}
Before the new one: say the Persian for courage, then the Persian for clever.
\end{quote}

\subsection*{You'll want to know: \fa{شوخ}}
\textbf{\fa{شوخ}} (\emph{shukh}) — \textquotedblleft{}witty, playful\textquotedblright{}.

\begin{quote}
\fa{آموزگار} \fa{ما} \fa{خیلی} \fa{شوخ} \fa{است} \fa{و} \fa{کلاس} \fa{را} \fa{شاد} \fa{میکند}.
\end{quote}

\begin{quote}
\emph{âmuzgâr-e mâ kheyli shukh ast va kelâs râ shâd mikonad.}
\end{quote}

\begin{quote}
Our teacher is very witty and cheers up the class.
\end{quote}

\emph{shukhi} 'a joke' and \emph{shukhi kardan} 'to joke' are built on it.

\begin{cousinweb}
\emph{shukh} is usually taken as a native Persian word, and \emph{shukhi}, 'a joke', is built on it. No source used here gives its older form, and its origin is not settled.
\end{cousinweb}

\begin{grammarlens}[title={What you've built}]
One more word for telling, explaining and describing at B1.
\end{grammarlens}

\subsection*{Your turn}
\begin{itemize}
\raggedright
  \item \emph{Say it:} \emph{shukh}
  \item \emph{Say it:} \emph{shukh}, once more
  \item \emph{Say it:} a sentence of your own with \emph{shukh}
  \item \emph{Recall:} say the Persian for courage, then the Persian for clever, then say \emph{shukh} again
\end{itemize}

\begin{tcolorbox}[breakable,colback=teal!4,colframe=teal!35!black,title={Before you move on}]
What does \fa{شوخ} mean? (\textquotedblleft{}Witty, playful\textquotedblright{}.) Say it once more.
\end{tcolorbox}
\section[\texorpdfstring{\fa{خوش} \fa{اخلاق} (khosh akhlâq)}{خوش اخلاق (khosh akhlâq)}]{\fa{خوش} \fa{اخلاق} (khosh akhlâq) — good-natured, pleasant}

\label{lesson:FA-C488-khosh-akhlaq}



\begin{quote}
Before the new one: say the Persian for clever, then the Persian for witty.
\end{quote}

\subsection*{You'll want to know: \fa{خوش} \fa{اخلاق}}
\textbf{\fa{خوش} \fa{اخلاق}} (\emph{khosh akhlâq}) — \textquotedblleft{}good-natured, pleasant\textquotedblright{}.

\begin{quote}
\fa{همسایه} \fa{جدید} \fa{ما} \fa{خیلی} \fa{خوش} \fa{اخلاق} \fa{است}.
\end{quote}

\begin{quote}
\emph{hamsâye-ye jadid-e mâ kheyli khosh akhlâq ast.}
\end{quote}

\begin{quote}
Our new neighbour is very pleasant.
\end{quote}

\emph{khosh} 'good' + \emph{akhlâq} 'manners, character'. The opposite is \emph{badakhlâq}.

\begin{cousinweb}
\emph{khosh akhlâq} is built in Persian on native \emph{khosh}, 'good', and \emph{akhlâq}, 'character, manners', which was borrowed from Arabic \emph{aḫlāq}.
\end{cousinweb}

\begin{grammarlens}[title={What you've built}]
One more word for telling, explaining and describing at B1.
\end{grammarlens}

\subsection*{Your turn}
\begin{itemize}
\raggedright
  \item \emph{Say it:} \emph{khosh akhlâq}
  \item \emph{Say it:} \emph{khosh akhlâq}, once more
  \item \emph{Say it:} a sentence of your own with \emph{khosh akhlâq}
  \item \emph{Recall:} say the Persian for clever, then the Persian for witty, then say \emph{khosh akhlâq} again
\end{itemize}

\begin{tcolorbox}[breakable,colback=teal!4,colframe=teal!35!black,title={Before you move on}]
What does \fa{خوش} \fa{اخلاق} mean? (\textquotedblleft{}Good-natured, pleasant\textquotedblright{}.) Say it once more.
\end{tcolorbox}
\section[\texorpdfstring{\fa{سخاوتمند} (sekhâvatmand)}{سخاوتمند (sekhâvatmand)}]{\fa{سخاوتمند} (sekhâvatmand) — generous}

\label{lesson:FA-C488-sekhavatmand}



\begin{quote}
Before the new one: say the Persian for witty, then the Persian for good-natured.
\end{quote}

\subsection*{You'll want to know: \fa{سخاوتمند}}
\textbf{\fa{سخاوتمند}} (\emph{sekhâvatmand}) — \textquotedblleft{}generous\textquotedblright{}.

\begin{quote}
\fa{او} \fa{مرد} \fa{سخاوتمندی} \fa{است} \fa{و} \fa{به} \fa{همه} \fa{کمک} \fa{میکند}.
\end{quote}

\begin{quote}
\emph{u mard-e sekhâvatmandi ast va be hame komak mikonad.}
\end{quote}

\begin{quote}
He is a generous man and helps everyone.
\end{quote}

\emph{sekhâvat} 'generosity' + Persian \emph{-mand} 'having'. In speech also \emph{dast-o-delbâz}.

\begin{cousinweb}
\emph{sekhâvatmand} is built in Persian on \emph{sekhâvat}, 'generosity', a loan usually traced to Arabic \emph{saḫāwa}, with the native ending \emph{-mand}, 'having'. EtymDB has no entry for it.
\end{cousinweb}

\begin{grammarlens}[title={What you've built}]
One more word for telling, explaining and describing at B1.
\end{grammarlens}

\subsection*{Your turn}
\begin{itemize}
\raggedright
  \item \emph{Say it:} \emph{sekhâvatmand}
  \item \emph{Say it:} \emph{sekhâvatmand}, once more
  \item \emph{Say it:} a sentence of your own with \emph{sekhâvatmand}
  \item \emph{Recall:} say the Persian for witty, then the Persian for good-natured, then say \emph{sekhâvatmand} again
\end{itemize}

\begin{tcolorbox}[breakable,colback=teal!4,colframe=teal!35!black,title={Before you move on}]
What does \fa{سخاوتمند} mean? (\textquotedblleft{}Generous\textquotedblright{}.) Say it once more.
\end{tcolorbox}
\section[\texorpdfstring{\fa{دلسوز} (delsuz)}{دلسوز (delsuz)}]{\fa{دلسوز} (delsuz) — caring, compassionate}

\label{lesson:FA-C488-delsuz}



\begin{quote}
Before the new one: say the Persian for good-natured, then the Persian for generous.
\end{quote}

\subsection*{You'll want to know: \fa{دلسوز}}
\textbf{\fa{دلسوز}} (\emph{delsuz}) — \textquotedblleft{}caring, compassionate\textquotedblright{}.

\begin{quote}
\fa{پرستارهای} \fa{این} \fa{بیمارستان} \fa{خیلی} \fa{دلسوز} \fa{هستند}.
\end{quote}

\begin{quote}
\emph{parastârhâ-ye in bimârestân kheyli delsuz hastand.}
\end{quote}

\begin{quote}
The nurses at this hospital are very caring.
\end{quote}

\emph{del} 'heart' + \emph{suz} 'burning' (stem of \emph{sukhtan}): a heart that burns for others.

\begin{cousinweb}
\emph{delsuz} is built on two native Persian pieces: \emph{del}, 'heart', and \emph{suz}, which EtymDB derives from \emph{sukhtan}, 'to burn' (EtymDB).
\end{cousinweb}

\begin{grammarlens}[title={What you've built}]
One more word for telling, explaining and describing at B1.
\end{grammarlens}

\subsection*{Your turn}
\begin{itemize}
\raggedright
  \item \emph{Say it:} \emph{delsuz}
  \item \emph{Say it:} \emph{delsuz}, once more
  \item \emph{Say it:} a sentence of your own with \emph{delsuz}
  \item \emph{Recall:} say the Persian for good-natured, then the Persian for generous, then say \emph{delsuz} again
\end{itemize}

\begin{tcolorbox}[breakable,colback=teal!4,colframe=teal!35!black,title={Before you move on}]
What does \fa{دلسوز} mean? (\textquotedblleft{}Caring, compassionate\textquotedblright{}.) Say it once more.
\end{tcolorbox}
